\documentclass[12pt]{article}
\usepackage[T1]{fontenc}
\usepackage[utf8]{inputenc}
\usepackage{lmodern}
\usepackage{textcomp}
\usepackage{enumitem}
\usepackage{geometry}
\geometry{margin=1in}
\begin{document}
\begin{center}
Answer Key - Health Sciences - K-12 Level Assessment 18 \
\small (Answers are ordered exactly as in the question paper.)
\end{center}

\section*{Multiple Choice}
\begin{enumerate}[label=\arabic*.]
\item \textbf{Answer:} D
\\ \textit{Options:} A) Always requires parenteral nutritional suport; B) May follow colectomy; C) Is always permanent; D) May benefit from sodium supplementation of \textgreater{}=90 mmol/l
\\ \textit{Note:} Intestinal failure can lead to malabsorption of nutrients and electrolytes, making sodium supplementation of 90 mmol/l or more beneficial to help maintain proper electrolyte balance and overall health.
\item \textbf{Answer:} C
\\ \textit{Options:} A) Protein-energy supplementation during pregnancy; B) Protein-energy supplementation during the first two years; C) Both of the interventions; D) None of the interventions
\\ \textit{Note:} Both interventions have been empirically demonstrated to enhance various aspects of child development, such as cognitive and physical growth, thereby supporting the overall health and well-being of children.
\item \textbf{Answer:} C
\\ \textit{Options:} A) Obesity; B) Sugar; C) Salt and salt-preserved foods; D) Dietary fibre
\\ \textit{Note:} Salt and salt-preserved foods are linked to an increased risk of stomach cancer because high salt intake can damage the stomach lining and promote the growth of certain bacteria that are associated with cancer development.
\item \textbf{Answer:} B
\\ \textit{Options:} A) 75 th and 2 nd BMI centile; B) 91 st and 2 nd BMI centile; C) 99.6 th and 0.4 th BMI centile; D) 75 th and 25 th BMI centile
\\ \textit{Note:} The correct answer is based on the understanding that a BMI between the 91 st and 2 nd centile indicates a range where a child's growth is considered normal compared to their peers, reflecting healthy development in weight relative to height.
\item \textbf{Answer:} D
\\ \textit{Options:} A) Religious belief; B) For better health; C) To reduce global carbon dioxide and methane emissions; D) None of the options given are correct
\\ \textit{Note:} The correct answer is "None of the options given are correct" because all the listed reasons for following a vegetarian diet, such as health benefits, ethical considerations, and environmental impact, are valid and commonly accepted motivations.
\end{enumerate}

\section*{True / False}
\begin{enumerate}[label=\arabic*.]
\setcounter{enumi}{5}
\item \textbf{Answer:} False
\\ \textit{Options:} A) True; B) False
\\ \textit{Note:} Cognitive decline in older adults is influenced by multiple factors, including genetics, lifestyle, and overall nutrition, rather than solely low dietary fat intake.
\item \textbf{Answer:} True
\\ \textit{Options:} A) True; B) False
\\ \textit{Note:} The elderly require more protein-dense foods because their bodies experience a decrease in muscle mass and protein synthesis, necessitating higher protein intake to maintain muscle health and overall nutritional balance.
\item \textbf{Answer:} False
\\ \textit{Options:} A) True; B) False
\\ \textit{Note:} Malnutrition in older adults often results in complications that can prolong hospital stays due to the need for additional medical interventions and recovery time.
\item \textbf{Answer:} True
\\ \textit{Options:} A) True; B) False
\\ \textit{Note:} mRNA, or messenger RNA, is indeed a type of coding RNA that transcribes genetic information from DNA and transports it to the ribosome, where it serves as a template for synthesizing proteins.
\item \textbf{Answer:} False
\\ \textit{Options:} A) True; B) False
\\ \textit{Note:} DNMT 3 a is primarily involved in de novo methylation, establishing new methylation marks rather than copying existing ones during DNA replication, as this process is primarily carried out by DNMT 1.
\end{enumerate}

\section*{Long Form Response}
\begin{enumerate}[label=\arabic*.]
\setcounter{enumi}{10}
\item \textbf{Answer:} Research indicates that respiratory diseases tend to be less common among vegetarians compared to non-vegetarians. This difference may be attributed to several factors, including the higher intake of fruits and vegetables among vegetarians, which provide essential nutrients and antioxidants that support lung health. Additionally, vegetarians often have lower rates of smoking and obesity, both of which are significant risk factors for respiratory issues. Furthermore, the anti-inflammatory properties of a plant-based diet may reduce the incidence of conditions like asthma and chronic bronchitis. Overall, the dietary choices and lifestyle habits associated with vegetarianism contribute to better respiratory health outcomes.
\\ \textit{Note:} correct because it identifies that the lower prevalence of respiratory diseases among vegetarians can be linked to their higher consumption of nutrient-rich foods, lower smoking rates, and reduced obesity, all of which contribute positively to lung health.
\item \textbf{Answer:} Vitamin K deficiency significantly impacts the blood clotting process because vitamin K is essential for synthesizing certain proteins known as clotting factors. Without adequate vitamin K, the body struggles to form clots effectively, leading to prolonged clotting times. This means that when a person sustains an injury, their blood takes longer to clot, increasing the risk of excessive bleeding. Prolonged clotting time can have serious health implications, such as increased risk of hemorrhagic events, which can be dangerous and may require medical intervention. Therefore, maintaining sufficient vitamin K levels is crucial for overall health and effective blood clotting.
\\ \textit{Note:} correct because it clearly articulates the crucial role of vitamin K in synthesizing clotting factors, explains how its deficiency leads to prolonged clotting times, and outlines the serious health risks associated with excessive bleeding due to delayed clot formation.
\end{enumerate}

\vfill
\noindent\textit{End of Answer Key}
\end{document}